\documentclass[12pt,spanish,letterpaper,oneside]{article}
\def\documenttitle {Las generaciones de los derechos humanos y sus antecedentes}
\def\documentsubtitle {Primera generación y Carta Magna}
\def\documentsubject {Primera generación de los derechos humanos y Carta Magna}
\def\documentauthor {Martin Jonathan de la Cruz}
\def\coursename {Antecedentes de los derechos humanos}
\def\coursecode {005}
\def\universityname {Universidad Abierta y a Distancia de México}
\def\universityfaculty {Licenciatura en Derecho}
\def\universitydepartment {Antecedentes de los derechos humanos}
\def\universitydepartmentimage {departamentos/UnADM}
\def\universitydepartmentimagecfg {height=1.57cm}
\def\universitylocation {Ciudad de México}
\def\authortable {\begin{tabular}{ll}
  \textbf{Alumno:} & Martin Jonathan de la Cruz \\
  Matricula: & ES2611202040 \\
  Figura docente: & Daniel Alexis Lozano Ortega \\
  Semestre/Bloque: & 2 / 2 \\
  Semana: & 3 \\
  Fecha: & 18 de octubre de 2026
\end{tabular}}
\usepackage{iftex}
\ifLuaTeX
  \let\pdfpagewidth\pagewidth
  \let\pdfpageheight\pageheight
\fi
\input{template}
\usepackage{fontspec}
\setmainfont{arial.ttf}[Path=C:/Windows/Fonts/,BoldFont=arialbd.ttf,ItalicFont=ariali.ttf,BoldItalicFont=arialbi.ttf]
\setsansfont{arial.ttf}[Path=C:/Windows/Fonts/,BoldFont=arialbd.ttf,ItalicFont=ariali.ttf,BoldItalicFont=arialbi.ttf]
\renewcommand{\familydefault}{\sfdefault}
\setcitestyle{authoryear,open={(},close={)}}
\usepackage{tikz}
\usepackage{eso-pic}
\newcommand{\insertcoverwatermark}{%
  \AddToShipoutPictureBG*{%
    \begin{tikzpicture}[remember picture,overlay]
      \node[opacity=0.16,inner sep=0pt] at (current page.center)
        {\includegraphics[width=0.72\paperwidth]{img/departamentos/UnADM.pdf}};
    \end{tikzpicture}%
  }%
}
\begin{document}
\insertcoverwatermark
\templatePortrait
\templatePagecfg
\templateFinalcfg
\newgeometry{margin=2.5cm}
\fontsize{12}{14.5}\selectfont
\onehalfspacing

\section{Introducción}
La clasificación por generaciones permite situar el reconocimiento de los derechos civiles y políticos en procesos históricos concretos, sin establecer una jerarquía respecto de otros derechos. Distinguir los límites medievales al poder del reconocimiento posterior de libertades individuales evita atribuir a una época el concepto jurídico de otra. El valor de los antecedentes depende, por ello, tanto de lo que anticiparon como de las personas a quienes dejaron fuera \citep{reyCantor2022S3,parlamentoCarta1215ADH}.

\section{Libertades, contexto histórico y límites del poder}
Rey Cantor, Rey Anaya y Rey Anaya distinguen derechos civiles y políticos, derechos económicos, sociales y culturales, y derechos de solidaridad. Proponen examinar en cada grupo quién es titular, quién está obligado y qué conducta se exige. La cronología se cruza con el espacio y con conflictos políticos, religiosos y sociales: no describe una aparición simultánea ni definitiva de todos los derechos. La libertad, la igualdad y la fraternidad ayudan a explicar sus respectivos contextos, pero los autores reconocen objeciones al término \emph{generaciones}: su uso no debe sugerir que una etapa reemplaza a las anteriores \citep[pp.~72--76]{reyCantor2022S3}.

Las libertades civiles y políticas se afirman en las revoluciones inglesas de 1688--1689, norteamericana de 1776 y francesa de 1789. Rey Cantor y coautoras precisan que no se puede hablar de derechos humanos en sentido estricto durante la Antigüedad o la Edad Media. Esta distinción impide proyectar la noción contemporánea de universalidad sobre instituciones anteriores \citep[p.~83]{reyCantor2022S3}.

La reconstrucción histórica no empieza ni termina con esas tres revoluciones. Los autores proponen observar tiempo y lugar, así como conflictos políticos y religiosos, para explicar por qué en determinadas circunstancias se exigieron nuevas libertades. En particular, vinculan los debates de tolerancia con las tensiones entre confesiones cristianas y con la Reforma protestante; sitúan las libertades individuales en el marco del Estado liberal, mientras que los derechos sociales se asocian con otro momento del Estado de derecho. Esta secuencia describe cambios en las demandas y en su reconocimiento, no la desaparición de derechos anteriores \citep[pp.~73--75]{reyCantor2022S3}. Por ello, la noción de primera generación debe leerse como un recurso para comprender las condiciones de afirmación de las libertades civiles y políticas, no como un criterio para negar el carácter interdependiente de los derechos actuales.

La Carta Magna de 1215 es un antecedente, no una declaración universal de derechos: resultó de un conflicto entre el rey Juan y los barones ingleses y puso límites escritos al ejercicio del poder regio. Sus garantías respondían a la estructura social de su época y no se extendían por igual a toda la población. Su relevancia histórica radica en la idea de que la autoridad puede quedar sometida a reglas y procedimientos, idea retomada después en otros debates constitucionales \citep{parlamentoCarta1215ADH}. Por ello conviene distinguir el antecedente medieval de la posterior formulación liberal de libertades civiles y políticas.

Reconocer un derecho, justificarlo y garantizarlo son operaciones distintas. La clasificación histórica orienta la explicación de cada reivindicación; valorar su protección actual exige identificar también las obligaciones jurídicas concretas y los medios para exigirlas. La dimensión histórica de los derechos no sustituye el análisis de sus garantías actuales.

\section{Conclusión}
El reconocimiento de libertades y garantías expresa una continuidad de luchas y transformaciones, no la sustitución de derechos antiguos por otros nuevos. La clasificación generacional aclara sus condiciones históricas si no se interpreta como una división absoluta. Es particularmente importante estudiar antecedentes como la Carta Magna sin idealizarlos: reconocer sus límites tempranos al poder y sus exclusiones permite valorar la igualdad jurídica posterior y exigir que la protección actual alcance a todas las personas.\footnote{Uso de IA: GitHub Copilot para organización conceptual y revisión lingüística; las afirmaciones históricas se contrastan con las fuentes citadas.}

\clearpage
\begin{thebibliography}{2}
\bibitem[Rey Cantor et~al.(2022)]{reyCantor2022S3}
Rey Cantor, E., Rey Anaya, G. A., \& Rey Anaya, A. M. (2022). \emph{Las generaciones de los derechos humanos: Libertad, igualdad y fraternidad} (10.\textsuperscript{a} ed.). Universidad Libre.
\bibitem[UK Parliament(s.~f.)]{parlamentoCarta1215ADH}
UK Parliament. (s.~f.). \emph{Magna Carta}. \url{https://www.parliament.uk/about/living-heritage/evolutionofparliament/originsofparliament/birthofparliament/overview/magnacarta/}
\end{thebibliography}
\end{document}